% =====================================================================
\chapter{Resolución de triángulos y producto escalar}
\label{tema:3}
% =====================================================================

Este tema extiende la trigonometría a triángulos cualesquiera mediante los teoremas del seno y del coseno, clasifica los casos de resolución de triángulos (incluido el caso ambiguo) y muestra que el teorema del coseno es la misma afirmación que la fórmula geométrica del producto escalar. Se sigue~\cite[§10.1, §10.2 y §10.8]{osat}, el tratamiento de los triángulos oblicuos de~\cite{stewart} y, para el producto escalar, la sección~1.2 de Strang~\cite{strang}.

\section{Área y teorema del seno}

En lo que sigue, $ABC$ es un triángulo con ángulos $A,B,C$ (que se identifican con sus medidas en radianes) y lados $a=BC$, $b=CA$, $c=AB$, de modo que cada lado es el opuesto al ángulo de la misma letra. Se cumple $A+B+C=\pi$ y todos los ángulos están en $(0,\pi)$, de modo que sus senos son positivos.

\begin{proposicion}[Área]\label{prop:area}
El área de $ABC$ es $S=\tfrac12 ab\sin C=\tfrac12 bc\sin A=\tfrac12 ca\sin B$.
\end{proposicion}

\begin{proof}
Se toma $CA$ como base, de longitud $b$. Sea $H$ el pie de la perpendicular desde $B$ a la recta $CA$. Si $C$ es agudo, el triángulo $BHC$ es rectángulo en $H$ con hipotenusa $a$ y ángulo $C$, luego la altura es $BH=a\sin C$. Si $C$ es obtuso, $H$ cae fuera del lado y el ángulo en $C$ del triángulo $BHC$ es $\pi-C$, luego $BH=a\sin(\pi-C)=a\sin C$ por el \cref{teo:simetrias}(iv). Si $C=\pi/2$, $BH=a=a\sin C$. En todos los casos $S=\tfrac12\,b\cdot a\sin C$; las otras dos expresiones se obtienen igual con otra base.
\end{proof}

\begin{teorema}[Teorema del seno]\label{teo:seno}
En todo triángulo,
\[
\frac{a}{\sin A}=\frac{b}{\sin B}=\frac{c}{\sin C}.
\]
\end{teorema}

\begin{proof}
Por la \cref{prop:area}, $\tfrac12 bc\sin A=\tfrac12 ca\sin B=\tfrac12 ab\sin C$. Dividiendo los tres miembros por $\tfrac12 abc>0$ se obtiene $\frac{\sin A}{a}=\frac{\sin B}{b}=\frac{\sin C}{c}$, y basta tomar inversos.
\end{proof}

El valor común de los tres cocientes es el diámetro $2R$ de la circunferencia circunscrita (\cref{eje:3-circunradio}).

\section{Teorema del coseno}

\begin{teorema}[Teorema del coseno]\label{teo:coseno}
En todo triángulo, $c^{2}=a^{2}+b^{2}-2ab\cos C$, y análogamente para los otros dos lados.
\end{teorema}

\begin{proof}
Se colocan $C$ en el origen, $B=(a,0)$ sobre el semieje positivo de abscisas y $A$ en el semiplano superior. Como $CA=b$ y el ángulo en $C$ mide $C$, la \cref{def:circunferencia} (escalada por $b$) da $A=(b\cos C,\,b\sin C)$, sea $C$ agudo, recto u obtuso. Entonces
\begin{align*}
c^{2}=\abs{AB}^{2}&=(b\cos C-a)^{2}+(b\sin C)^{2}\\
&=b^{2}(\cos^{2}C+\sin^{2}C)-2ab\cos C+a^{2}=a^{2}+b^{2}-2ab\cos C.\qedhere
\end{align*}
\end{proof}

\begin{corolario}\label{cor:pitagoras}
$C=\pi/2$ si y solo si $c^{2}=a^{2}+b^{2}$ (teorema de Pitágoras y su recíproco); $C$ es obtuso si y solo si $c^{2}>a^{2}+b^{2}$, y agudo si y solo si $c^{2}<a^{2}+b^{2}$.
\end{corolario}

\begin{proof}
Por el \cref{teo:coseno}, $a^{2}+b^{2}-c^{2}=2ab\cos C$, que tiene el signo de $\cos C$. En $(0,\pi)$, $\cos C$ es positivo, nulo o negativo según que $C$ sea agudo, recto u obtuso (\cref{cor:signos}).
\end{proof}

\section{Resolución de triángulos}

\emph{Resolver} un triángulo es calcular sus seis elementos a partir de tres datos que incluyan al menos un lado. La \cref{tab:casos} resume qué teorema se usa en cada caso~\cite[§10.1--10.2]{osat}.

\begin{table}[ht]
\centering\small
\caption{Casos de resolución de triángulos.}\label{tab:casos}
\begin{tabular}{@{}>{\raggedright\arraybackslash}p{4.6cm}>{\raggedright\arraybackslash}p{4.6cm}>{\raggedright\arraybackslash}p{4.6cm}@{}}
\toprule
Datos & Primer paso & Número de soluciones \\
\midrule
LLL (tres lados) & teorema del coseno para un ángulo & una, si se cumple la desigualdad triangular \\
LAL (dos lados y el ángulo comprendido) & teorema del coseno para el tercer lado & una \\
ALA o AAL (dos ángulos y un lado) & $C=\pi-A-B$ y teorema del seno & una, si $A+B<\pi$ \\
LLA (dos lados y un ángulo no comprendido) & teorema del seno & cero, una o dos (caso ambiguo) \\
\bottomrule
\end{tabular}
\end{table}

\begin{proposicion}[Caso ambiguo]\label{prop:ambiguo}
Se dan $a$, $b$ y el ángulo $A$ opuesto a $a$, con $A$ agudo, y sea $h=b\sin A$. Entonces existen: ningún triángulo si $a<h$; uno (rectángulo en $B$) si $a=h$; dos si $h<a<b$; y uno si $a\ge b$.
\end{proposicion}

\begin{proof}
Por el \cref{teo:seno}, $\sin B=b\sin A/a=h/a$. Si $a<h$, sería $\sin B>1$, imposible. Si $a=h$, $\sin B=1$ y $B=\pi/2$. Si $a>h$, hay dos ángulos en $(0,\pi)$ con ese seno, $B_1\in(0,\pi/2)$ y $B_2=\pi-B_1$, y cada uno da un triángulo si y solo si $A+B<\pi$. Para $B_1$ siempre se cumple, porque $A$ y $B_1$ son agudos. Para $B_2$, la condición $A+\pi-B_1<\pi$ equivale a $A<B_1$, es decir (seno creciente en $[0,\pi/2]$) a $\sin A<\sin B_1=b\sin A/a$, o sea, a $a<b$.
\end{proof}

\begin{ejemplo}[Caso LAL]\label{ej:lal}
Sean $a=8$, $b=5$ y $C=\pi/3$. Por el \cref{teo:coseno}, $c^{2}=64+25-2\cdot 8\cdot 5\cdot\tfrac12=49$, luego $c=7$. Por el mismo teorema, $\cos A=\frac{b^{2}+c^{2}-a^{2}}{2bc}=\frac{25+49-64}{70}=\frac17$, de donde $A\approx 81{,}79^{\circ}$, y $B=180^{\circ}-60^{\circ}-81{,}79^{\circ}=38{,}21^{\circ}$. El área es $\tfrac12\cdot 8\cdot 5\cdot\frac{\sqrt3}{2}=10\sqrt3\approx 17{,}32$.
\end{ejemplo}

Cuando se conocen los tres lados, el área puede calcularse sin ángulos mediante la \emph{fórmula de Herón}, $S=\sqrt{s(s-a)(s-b)(s-c)}$ con $s=\frac{a+b+c}{2}$, que se demuestra en el \cref{eje:3-heron}.

\section{Producto escalar y ángulo entre vectores}

\begin{definicion}[Producto escalar]
Para $u=(u_1,\dots,u_n)$ y $v=(v_1,\dots,v_n)$ en $\R^{n}$, el \emph{producto escalar} es $u\cdot v=\sum_{i=1}^{n}u_iv_i$ y la \emph{norma} es $\norm{u}=\sqrt{u\cdot u}$.
\end{definicion}

El producto escalar es simétrico y bilineal; en particular,
\begin{equation}\label{eq:norma-diferencia}
\norm{u-v}^{2}=(u-v)\cdot(u-v)=\norm{u}^{2}+\norm{v}^{2}-2\,u\cdot v .
\end{equation}

\begin{teorema}[Forma geométrica del producto escalar]\label{teo:escalar}
Sean $u,v\in\R^{2}$ (o $\R^{3}$) no nulos y $\varphi\in[0,\pi]$ el ángulo que forman. Entonces
\[
u\cdot v=\norm{u}\,\norm{v}\cos\varphi .
\]
\end{teorema}

\begin{proof}
Si $u$ y $v$ no son paralelos, los puntos $O$, $u$ y $v$ forman un triángulo con lados $\norm{u}$, $\norm{v}$ y $\norm{u-v}$, en el que el ángulo opuesto al lado $\norm{u-v}$ es $\varphi$. El \cref{teo:coseno} da $\norm{u-v}^{2}=\norm{u}^{2}+\norm{v}^{2}-2\norm{u}\norm{v}\cos\varphi$, y comparando con~\eqref{eq:norma-diferencia} se obtiene el enunciado. Si son paralelos, $v=\lambda u$ con $\lambda\ne0$: si $\lambda>0$, $\varphi=0$ y $u\cdot v=\lambda\norm{u}^{2}=\norm{u}\norm{v}$; si $\lambda<0$, $\varphi=\pi$ y $u\cdot v=\lambda\norm{u}^{2}=-\norm{u}\norm{v}$.
\end{proof}

\begin{corolario}\label{cor:escalar}
Para $u,v$ no nulos del plano o del espacio: (i) $u\perp v$ si y solo si $u\cdot v=0$; (ii) $\abs{u\cdot v}\le\norm{u}\norm{v}$ (desigualdad de Cauchy--Schwarz); (iii) la componente de $u$ en la dirección de $v$ es $\norm{u}\cos\varphi=u\cdot v/\norm{v}$, y la proyección de $u$ sobre la recta de $v$ es $\frac{u\cdot v}{\norm{v}^{2}}\,v$.
\end{corolario}

\begin{proof}
(i) $\cos\varphi=0$ si y solo si $\varphi=\pi/2$ en $[0,\pi]$. (ii) Se sigue de $\abs{\cos\varphi}\le1$. (iii) En el triángulo rectángulo formado por $u$ y su proyección, el cateto contiguo a $\varphi$ mide $\norm{u}\cos\varphi$ (con signo negativo si $\varphi>\pi/2$); se multiplica por el vector unitario $v/\norm{v}$.
\end{proof}

En $\R^{n}$ con $n>3$ no hay triángulos que dibujar y se procede al revés: la desigualdad de Cauchy--Schwarz se demuestra algebraicamente y se \emph{define} el ángulo por $\cos\varphi=u\cdot v/(\norm{u}\norm{v})$ (módulo~2, tema~2). El cociente recibe el nombre de \emph{similitud coseno}.

\begin{ejemplo}\label{ej:angulo}
Para $u=(1,1,0)$ y $v=(1,0,1)$: $u\cdot v=1$, $\norm{u}=\norm{v}=\sqrt2$, luego $\cos\varphi=1/2$ y $\varphi=\pi/3$. Los tres puntos $O$, $u$, $v$ forman un triángulo equilátero de lado $\sqrt2$, ya que $\norm{u-v}=\norm{(0,1,-1)}=\sqrt2$.
\end{ejemplo}

\begin{observacion}[Errores frecuentes]
(a) En el caso LLA, olvidar la segunda solución $B_2=\pi-B_1$, porque la calculadora solo devuelve $\arcsin$ en $[-\pi/2,\pi/2]$. (b) Usar el teorema del seno para obtener el ángulo mayor de un triángulo: si es obtuso, el arcoseno da su suplementario. Es más seguro obtener primero el ángulo menor, o usar el teorema del coseno, que distingue agudos de obtusos por el signo. (c) Confundir el ángulo entre vectores con el ángulo entre rectas: el primero está en $[0,\pi]$.
\end{observacion}

\begin{observacion}[Conexión con el curso]
La demostración del \cref{teo:escalar} es la de la proposición que relaciona el ángulo geométrico con el producto escalar en el módulo~2 (tema~2). En el Transformer, la puntuación de atención entre una consulta $q$ y una clave $k$ es $q\cdot k/\sqrt{d_k}=\norm{q}\norm{k}\cos\varphi/\sqrt{d_k}$ (módulo~13): crece con las normas y con la alineación de ambas direcciones. RoPE no cambia las normas y modifica solo el ángulo $\varphi$ en función de la distancia entre posiciones (módulo~14, tema~4).
\end{observacion}

\begin{lecturas}
\item OpenStax, \emph{Algebra and Trigonometry}~\cite{osat}, §10.1 «Non-right Triangles: Law of Sines», §10.2 «Non-right Triangles: Law of Cosines» y §10.8 «Vectors».
\item Strang~\cite{strang}, §1.2 «Lengths and Dot Products».
\end{lecturas}

\section*{Ejercicios del tema 3}
\addcontentsline{toc}{section}{Ejercicios del tema 3}

\begin{ejercicio}{1}{eje:3-que-teorema}
Indique qué teorema usaría en primer lugar y cuántas soluciones puede haber en cada caso: (a) $a=4$, $b=6$, $c=9$; (b) $A=40^{\circ}$, $B=75^{\circ}$, $c=12$; (c) $b=7$, $c=10$, $A=110^{\circ}$; (d) $a=5$, $b=8$, $A=30^{\circ}$.
\end{ejercicio}

\begin{ejercicio}{1}{eje:3-imposible}
¿Existe un triángulo con lados $3$, $4$ y $8$? ¿Y con $a=2$, $b=5$ y $A=60^{\circ}$? Justifique ambas respuestas.
\end{ejercicio}

\begin{ejercicio}{1}{eje:3-obtuso}
Sin calcular ángulos, decida si el triángulo de lados $6$, $7$ y $10$ es acutángulo, rectángulo u obtusángulo.
\end{ejercicio}

\begin{ejercicio}{2}{eje:3-ala}
Resuelva el triángulo con $A=45^{\circ}$, $B=60^{\circ}$ y $c=10$, y calcule su área.
\fuente{adaptado de OpenStax, \emph{Algebra and Trigonometry}~\cite{osat}, §10.1, ejercicios de la sección sobre los casos ALA y AAL.}
\end{ejercicio}

\begin{ejercicio}{2}{eje:3-ambiguo}
Sean $a=6$, $b=8$ y $A=35^{\circ}$. Determine cuántos triángulos existen y resuélvalos todos. Use $\sin35^{\circ}\approx0{,}5736$.
\fuente{adaptado de OpenStax, \emph{Algebra and Trigonometry}~\cite{osat}, §10.1, ejercicios de la sección sobre el caso ambiguo.}
\end{ejercicio}

\begin{ejercicio}{2}{eje:3-lal}
Resuelva el triángulo con $b=5$, $c=7$ y $A=60^{\circ}$.
\fuente{adaptado de OpenStax, \emph{Algebra and Trigonometry}~\cite{osat}, §10.2, ejercicios de la sección sobre el caso LAL.}
\end{ejercicio}

\begin{ejercicio}{2}{eje:3-barcos}
Dos barcos salen a la vez de un puerto en direcciones que forman $120^{\circ}$. Cuando el primero ha recorrido 3 km y el segundo 5 km, ¿qué distancia los separa? ¿Qué ángulo forma, en el segundo barco, la visual al puerto con la visual al primer barco?
\fuente{adaptado de OpenStax, \emph{Algebra and Trigonometry}~\cite{osat}, §10.2, problemas de aplicación de navegación.}
\end{ejercicio}

\begin{ejercicio}{2}{eje:3-lll}
Resuelva el triángulo de lados $a=4$, $b=5$ y $c=6$, compruebe que $C=2A$ y calcule su área de dos maneras: con la \cref{prop:area} y con la fórmula de Herón.
\end{ejercicio}

\begin{ejercicio}{2}{eje:3-angulo-vectores}
Calcule el ángulo entre $u=(1,2,2)$ y $v=(2,-1,2)$, la componente de $u$ en la dirección de $v$ y la proyección de $u$ sobre la recta de $v$.
\fuente{adaptado de Strang~\cite{strang}, §1.2, problemas sobre ángulos entre vectores.}
\end{ejercicio}

\begin{ejercicio}{2}{eje:3-coseno-normas}
Se sabe que $\norm{u}=3$, $\norm{v}=4$ y $\norm{u-v}=\sqrt{13}$. Calcule $u\cdot v$, el ángulo entre $u$ y $v$ y $\norm{u+v}$.
\end{ejercicio}

\begin{ejercicio}{2}{eje:3-similitud}
Tres embeddings de dimensión~4 son $a=(1,0,1,1)$, $b=(0,1,1,1)$ y $c=(2,0,2,2)$. Calcule las similitudes coseno $\cos(a,b)$, $\cos(a,c)$ y $\cos(b,c)$ y los ángulos correspondientes. ¿Qué indica que $\cos(a,c)=1$ aunque $a\ne c$?
\end{ejercicio}

\begin{ejercicio}{3}{eje:3-circunradio}
Sea $R$ el radio de la circunferencia circunscrita al triángulo $ABC$. Demuestre que $\frac{a}{\sin A}=2R$. Indicación: trace el diámetro que pasa por $B$ y use el teorema del ángulo inscrito; distinga los casos $A$ agudo, recto y obtuso.
\end{ejercicio}

\begin{ejercicio}{3}{eje:3-heron}
Demuestre la fórmula de Herón $S=\sqrt{s(s-a)(s-b)(s-c)}$, con $s=\frac{a+b+c}{2}$, a partir de la \cref{prop:area} y del \cref{teo:coseno}. Indicación: calcule $16S^{2}=4a^{2}b^{2}(1-\cos^{2}C)$ y factorice como diferencia de cuadrados.
\end{ejercicio}

\begin{ejercicio}{3}{eje:3-proyecciones}
(a) Demuestre la \emph{ley de las proyecciones}: en todo triángulo $a=b\cos C+c\cos B$. (b) Escribiendo las tres leyes análogas y combinándolas, deduzca de nuevo el \cref{teo:coseno}.
\end{ejercicio}

\begin{ejercicio}{3}{eje:3-mayor-lado}
Demuestre que en todo triángulo $a>b$ si y solo si $A>B$: a mayor lado se opone mayor ángulo. Indicación: use el \cref{teo:seno} y el \cref{cor:suma-producto} para el signo de $\sin A-\sin B$.
\end{ejercicio}
